\textbf{Keywords:} Course Goals, Roles, Skills, Zen of IC Design, Design
Process, Tapeout

\section{Who}\label{who}

My name is

Carsten Wulff
\href{mailto:carstenw@ntnu.no}{\nolinkurl{carstenw@ntnu.no}}

I finished my Masters in 2002, and did a Ph.D on analog-to-digital
converters finished in 2008.

Since that time, I've had a three axis in my work/hobby life.

I work at \href{https://www.nordicsemi.com}{Nordic Semiconductor} where
I've been since 2008. The first 7 years I did analog design (ADCs,
DC/DCs, GPIO). The next 7 years I was the Wireless Group Manager. The
Wireless group made most of the analog and RF designs for Nordic's
short-range products. Now I'm the IC Scientist, and focus on technical
issues with our integrated circuits that occur before we go into volume
production.

I work at \href{https://ntnu.no}{NTNU} where I did a part time postdoc
from 2014 - 2017. From 2020 I've been working on and teaching
\href{https://www.ntnu.edu/studies/courses/TFE4188\#tab=omEmnet}{Advanced
Integrated Circuits}

I have a hobby trying to figure out how to make a new analog circuit
design paradigm. The one we have today with
schematic/simulation/layout/verification/simulation is too slow


{
\centering
\aicfig{media/timeline_tikz.png}
}

\emph{Figure 1: My life }

\section{How I see our roles}\label{how-i-see-our-roles}\index{how i see our roles@How I see our roles}

In Figure 2 you can see how I think about the research universe. There
are things we know to be possible, things that actually are impossible
(travel back in time, breaking thermodynamics, travel with a speed
beyond light).

Between the impossible, and the possible, lies the unknown. I consider
our roles as follows:

\textbf{Professors:} Guide students on what is impossible, possible, and
hints on what might be possible

\textbf{Ph.D students:} Venture into the unknown and make something
(more) possible

\textbf{Master students:} Learn all that is currently possible

\textbf{Bachelor students:} Learn how to make complicated into easy

\textbf{Industry:} Take what is possible, and/or complicated, and make
it easy

Everyone in the figure pushes a border: teaching moves knowledge from
complicated towards easy, research moves the border of the unknown, and
industry earns its living on the innermost rings. This course lives one
ring out from the centre - what is known and complicated, made
learnable.


{
\centering
\aicfig{media/l01_universe_tikz.png}
}

\emph{Figure 2: The research universe: Easy at the centre, then
Complicated, Possible and Unknown, with the Impossible beyond the last
border}

\section{I want you to learn the skills necessary to make your own
ICs}\label{i-want-you-to-learn-the-skills-necessary-to-make-your-own-ics}

In 2020 the global integrated circuit market was
\href{https://www.fortunebusinessinsights.com/integrated-circuit-market-106522}{437.7
billion dollars}! The market is expected to grow to 1136 billion in
2028. Integrated circuits enable all technologies.

I will be dead in approximately 50 years, and will retire in
approximately 20 years. Everything I know will be gone (except for the
small pieces I've left behind in videos or written word)

Someone must take over, and to do that, they need to know most of what I
know, and hopefully a bit more.

That's where some of you come in. Some of you will find integrated
circuits interesting to make, and in addition, you have the stamina,
patience, and brain necessary to learn some of the hardest topics in the
world.

\section{There will always be analog circuits, because the real world is
analog}\label{there-will-always-be-analog-circuits-because-the-real-world-is-analog}

In this course, we'll focus on analog ICs, because the real world is
analog, and all ICs must have some analog components, otherwise they
won't work.

The steps to make integrated circuits are split in two. We have an
analog flow, and a digital flow, as shown in Figure 3.

It's rare to find a single human that does both flows well. Usually
people choose, and I think it's based on what they like and their
personality.

If you like the world to be ordered, with definite answers, then it's
likely that you'll find the digital flow interesting.

If you're comfortable with not knowing, and have an insatiable desire to
understand how the world \emph{really} works at a fundamental level,
then it's likely that you'll find the analog flow interesting.


{
\centering
\aicfig{media/dig_des_tikz.png}
}

\emph{Figure 3: Analog and Digital design process }

\section{Will you tape-out an IC?}\label{will-you-tape-out-an-ic}\index{will you tape-out an ic?@Will you tape-out an IC?}

Something that would make me really happy is if someone is able to
tapeout an IC in this course.

It's now possible without signing an NDA or buying expensive software
licenses.

In 2020 Google and Skywater joined forces to release a 130 nm process
design kit to the public. In addition, they have fueled a renaissance of
open source software tools.

\href{https://tinytapeout.com}{tinytapeout} runs cheap shuttles which
makes it possible for a private citizen to tape-out their own integrated
circuit.

\subsection{What the team needs to know to design
ICs}\label{what-the-team-needs-to-know-to-design-ics}

There are a multitude of tools and skills needed to design professional
ICs. It's not likely that you'll find all the skills in one human, and
even if you could, one human does not have sufficient bandwidth to
design ICs with all its aspects in a reasonable timeline

That is, unless we can find a way to make ICs easier.

The skills needed are

\begin{itemize}
\tightlist
\item
  \emph{Project flow support}: \textbf{Confluence}, JIRA, risk
  management (DFMEA), failure analysis (8D)
\item
  \emph{Language}: \textbf{English}, \textbf{Writing English (Latex,
  Word, Email)}
\item
  \emph{Psychology}: Personalities, convincing people, presentations
  (Powerpoint, Deckset), \textbf{stress management (what makes your
  brain turn off?)}
\item
  \emph{DevOps}: \textbf{Linux}, build systems (CMake, make, ninja),
  continuous integration (bamboo, jenkins), \textbf{version control
  (git)}, containers (docker), container orchestration (swarm,
  kubernetes)
\item
  \emph{Programming}: Python, C, C++, Matlab \emph{Since 1999 I've
  programmed in Python, Go, Visual BASIC, PHP, Ruby, Perl, C\#, SKILL,
  Ocean, Verilog-A, C++, BASH, AWK, VHDL, SPICE, MATLAB, ASP, Java, C,
  SystemC, Verilog, Assembler, and probably a few I've forgotten.}
\item
  \emph{Firmware}: signal processing, algorithms, software architecture,
  security
\item
  \emph{Infrastructure}: \textbf{Power management}, \textbf{reset},
  \textbf{bias}, \textbf{clocks}
\item
  \emph{Domains}: CPUs, peripherals, memories, bus systems
\item
  \emph{Sub-systems}: \textbf{Radio's}, \textbf{analog-to-digital
  converters}, \textbf{comparators}
\item
  \emph{Blocks}: \textbf{Analog Radio}, Digital radio baseband
\item
  \emph{Modules}: Transmitter, \textbf{receiver}, de-modulator, timing
  recovery, state machines
\item
  \emph{Designs}: \textbf{Opamps}, \textbf{amplifiers},
  \textbf{current-mirrors}, adders, random access memory blocks,
  standard cells
\item
  \emph{Tools}: \textbf{schematic}, \textbf{layout}, \textbf{parasitic
  extraction}, synthesis, place-and-route, \textbf{simulation},
  (System)Verilog, \textbf{netlist}
\item
  \emph{Physics}: transistor, pn junctions, quantum mechanics
\end{itemize}

\subsection{Zen of IC design (stolen from Zen of
Python)}\label{zen-of-ic-design-stolen-from-zen-of-python}

When you learn something new, it's good to listen to someone that has
done whatever it is before.

Here are some guiding principles that you'll likely forget.

\begin{itemize}
\item
  Beautiful is better than ugly.
\item
  Explicit is better than implicit.
\item
  Simple is better than complex.
\item
  Complex is better than complicated.
\item
  Readability counts (especially schematics).
\item
  Special cases aren't special enough to break the rules.
\item
  Although practicality beats purity.
\item
  In the face of ambiguity, refuse the temptation to guess.
\item
  There should be one \textbf{and preferably only one} obvious way to do
  it.
\item
  Now is better than never.
\item
  Although never is often better than \emph{right} now.
\item
  If the implementation is hard to explain, it's a bad idea.
\item
  If the implementation is easy to explain, it may be a good idea.
\end{itemize}

\subsection{IC design mantra}\label{ic-design-mantra}\index{ic design mantra@IC design mantra}

To copy an old mantra I have on learning programming (run it in a
bash/zsh/cshrc terminal, or in your brain)

\begin{Shaded}
\begin{Highlighting}[]
\NormalTok{echo }\OtherTok{"}\StringTok{Find a problem that you really want to solve,}\OtherTok{"}\NormalTok{\textbackslash{}}
     \OtherTok{"}\StringTok{and learn programming to solve it.}\OtherTok{"}\NormalTok{\textbackslash{}}
     \OtherTok{"}\StringTok{There is no point in saying \textquotesingle{}I want to learn programming\textquotesingle{},}\OtherTok{"}\NormalTok{\textbackslash{}}
     \OtherTok{"}\StringTok{then sit down with a book to read about programming,}\OtherTok{"}\NormalTok{\textbackslash{}}
     \OtherTok{"}\StringTok{and expect that you will learn programming that way.}\OtherTok{"}\NormalTok{\textbackslash{}}
     \OtherTok{"}\StringTok{It will not happen. The only way to learn programming}\OtherTok{"}\NormalTok{\textbackslash{}}
     \OtherTok{"}\StringTok{is to do it, a lot.}\OtherTok{"}\NormalTok{ \textbackslash{}}
\NormalTok{     |perl {-}pe }\OtherTok{\textquotesingle{}}\SpecialStringTok{s/programming/analog design/ig}\OtherTok{\textquotesingle{}}
\end{Highlighting}
\end{Shaded}

\subsection{Analog Design Process}\label{analog-design-process}\index{analog design process@Analog design process}

\begin{itemize}
\item
  Define the problem, what are you trying to solve?
\item
  Find a circuit that can solve the problem (papers, books)
\item
  Find right transistor sizes. What transistors should be weak
  inversion, strong inversion, or don't care?
\item
  Write a verification plan (ask chat). Plan to simulate everything that
  could go wrong.
\item
  Check operating region of transistors (.op)
\item
  Check key parameters (.dc, .ac, .tran)
\item
  Check function. Exercise all inputs. Check all control signals
\item
  Check key parameters in all corners. Check mismatch (Monte-Carlo
  simulation)
\item
  Do layout, and check it's error free. Run design rule checks (DRC).
  Check layout versus schematic (LVS)
\item
  Extract parasitics from layout. Resistance, capacitance, and
  inductance if necessary.
\item
  On extracted parasitic netlist, check key parameters in all corners
  and mismatch (if possible).
\item
  If everything works, then you're done.
\end{itemize}

\emph{On failure, go back as far as necessary}

\section{My Goal}\label{my-goal}

Don't expect that I'll magically take information and put it inside your
head, and you'll suddenly understand everything about making ICs.

\textbf{You are the one that must teach yourself everything.}

I consider my role as a guide, similar to a mountain guide. I can't
carry you up the mountain, you need to walk up the mountain , but I know
the safe path to take and increase the likelihood that you'll come back
alive.

My guide role:

\begin{itemize}
\tightlist
\item
  Enable you to read the books on integrated circuits
\item
  Enable you to read papers (latest research)
\item
  Correct misunderstandings on the topic
\item
  Answer any questions you have on the chapters
\end{itemize}

I'm not a mind reader, I can't see inside your head. That means, you
must ask questions. Only by your questions can I start to understand
what pieces of information is missing from your head, or maybe somehow
correct your understanding.

At the same time, and similar to a mountain guide, you should not assume
I'm always right. I'm human, and I will make mistakes. And maybe you can
correct my understanding of something. All I care about is to
\emph{really} understand how the world works, so if you think my
understanding is wrong, then I'll happily discuss.

\section{Syllabus}\label{syllabus}

The syllabus will be from Analog Integrated Circuit Design by Carusone,
Johns and Martin {[}1{]}, which everyone calls CJM, and Circuits for all
seasons.

These lecture notes are a supplement to the book. I try to give some
background, and how to think about electronics. It's not my goal to
repeat information that you can find in the book.

Buy a hard-copy of the book if you don't have that. Don't expect to
understand the book by reading the PDF.

\section{Software}\label{software}

We'll use professional-grade open source software for everything: xschem
for schematics, ngspice for simulation, the SKY130A PDK, Magic VLSI and
netgen for layout and verification, and surfer, iverilog and verilator
on the digital side. No licence server stands between you and your
design.

Open source software (xschem, ngspice, sky130A PDK, Magic VLSI, netgen,
surfer, iverilog, verilator)

I've made a rather detailed (at least I think so myself) tutorial on how
to make a current mirror with the open source tools. You have to do that
tutorial. It's Milestone 0 of the project and does count towards your
final grade.

\href{https://analogicus.com/aic2026/sky130nm_tutorial}{Skywater 130 nm
Tutorial}

I've also made some more complex examples, that can be found at the link
below. There are digital logic cells, standard transistors, and few
other blocks.

\href{https://wulffern.github.io/aicex}{aicex}

\section{Summary}\label{summary}

The one-page version of this chapter:

\begin{itemize}
\tightlist
\item
  The real world is analog, so every IC carries analog circuits at its
  edges - there is no purely digital chip
\item
  Making an IC splits into an analog and a digital flow, and it is rare
  to find one human who does both well
\item
  Digital designers reuse each other's work; analog designers redraw -
  closing that gap is a theme of this course
\item
  What you need from here: the refresher chapters, a working toolchain,
  and the habit of simulating everything
\end{itemize}

\phantomsection\label{refs}
\begin{CSLReferences}{0}{0}
\bibitem[\citeproctext]{ref-cjm11}
\CSLLeftMargin{{[}1{]} }%
\CSLRightInline{T. C. Carusone, D. Johns, and K. Martin, \emph{Analog
integrated circuit design}. Wiley, 2011 {[}Online{]}. Available:
\url{https://books.google.no/books?id=1OIJZzLvVhcC}}

\end{CSLReferences}
